% Einheit 3 von 4 – Klammern auflösen (bank/terme/e3.jsonl, zusammenbau v0.1)
%% TODO Zweigzeile Teil 1: was hier gelernt wird (Satz in Schülersprache aus dem Einheitstitel) – Einheit 3
\einheitenkopf[e3]{Einheit 3 von 4 · Klammern auflösen}
\zweigzeile{ab Kl. 7, je nach Buch bis Kl. 8 · keine P10-Aufgabe}
\setcounter{aufgabe}{16}

%% TODO Titel: Ich-kann-Satz fehlt in der Bank (Platzhalter: Kettenname) – Nr. 17
%% TODO Anweisung: ein Satz über der ersten Teilaufgabe fehlt in der Bank – Nr. 17
\begin{aufgabe}{Zeichen vor der Klammer feststellen}
\begin{teile}
\teil $7 + (a - 2)$ – was steht vor der Klammer? \leerfeld
\end{teile}
\end{aufgabe}

%% TODO Titel: Ich-kann-Satz fehlt in der Bank (Platzhalter: Kettenname) – Nr. 18
%% TODO Anweisung: ein Satz über der ersten Teilaufgabe fehlt in der Bank – Nr. 18
\begin{aufgabe}{Klammern}
%% TODO Darstellung für schwach fehlt in der Bank (terme-e3-k2-s1-v1; Abschnitt „Für schwache Schüler“)
\swz{$\text{Löse die Klammer auf: $4x + (3y - 7)$}$}{}
%% TODO Darstellung für schwach fehlt in der Bank (terme-e3-k2-s1-v2; Abschnitt „Für schwache Schüler“)
\swz{$\text{Löse die Klammer auf: $2a + (5b + 1)$}$}{}
%% TODO Darstellung für schwach fehlt in der Bank (terme-e3-k2-s1-v3; Abschnitt „Für schwache Schüler“)
\swz{$\text{Löse die Klammer auf: $8 + (6x - y)$}$}{}
%% TODO Darstellung für schwach fehlt in der Bank (terme-e3-k2-s1-v4; Abschnitt „Für schwache Schüler“)
\swz{$\text{Löse die Klammer auf: $7y + (2z + 9)$}$}{}
%% TODO Darstellung für schwach fehlt in der Bank (terme-e3-k2-s1-v5; Abschnitt „Für schwache Schüler“)
\swz{$\text{Löse die Klammer auf: $3b + (4c - 5)$}$}{}
\swfrage{Was bleibt gleich, was ändert sich?}
\swa{Löse die Klammer mit der Tabelle auf: $5 \cdot (x + 3)$ \leerfeld}{\sachtabelle{ccc}{$\cdot$ & $x$ & $+3$}{$5$ & \leerzelle & \leerzelle}}
%% TODO Darstellung für schwach fehlt in der Bank (terme-e3-k2-s3-v1; Abschnitt „Für schwache Schüler“)
\swz{$\text{Löse die Klammer auf: $5a - (2b + 4)$}$}{}
%% TODO Darstellung für schwach fehlt in der Bank (terme-e3-k2-s4-v1; Abschnitt „Für schwache Schüler“)
\swz{$\text{Löse die Klammer auf: $7a - (2b - 5 + c)$}$}{}
%% TODO Darstellung für schwach fehlt in der Bank (terme-e3-k2-s5-v1; Abschnitt „Für schwache Schüler“)
\swz{$\text{Löse die Klammer auf: $-2 \cdot (x + 5)$}$}{}
%% TODO Darstellung für schwach fehlt in der Bank (terme-e3-k2-s6-v1; Abschnitt „Für schwache Schüler“)
\swz{$\text{Löse die Klammer auf: $5x + 2 \cdot (x + 6)$}$}{}
\swz[3]{$\text{Löse die Klammer auf: $3 \cdot (2x + 5) - 2 \cdot (x - 4)$}$}{}
\end{aufgabe}

%% TODO Titel: Ich-kann-Satz fehlt in der Bank (Platzhalter: Kettenname) – Nr. 19
%% TODO Anweisung: ein Satz über der ersten Teilaufgabe fehlt in der Bank – Nr. 19
\begin{aufgabe}{Klammern – Fehler finden, Begründen, Darstellungswechsel}
\swz[3]{Finde den Fehler und rechne richtig: \rechnung{12 - (a - 5) &= 12 - a - 5 \\ &= 7 - a}}{}
\swfrage{Sind $5 - (x - 2)$ und $7 - x$ gleichwertig? Begründe.}
\swz{Flächeninhalt des Rechtecks (Maße in cm) – Term mit Klammer und ohne Klammer?}{\viereck[seiten={x+2,5,x+2,5}]{(0,0)}{(6,0)}{(6,2.5)}{(0,2.5)}}
\end{aufgabe}

\uebersichtskasten{Plus vor der Klammer: Klammer weglassen. \quad a + (b − c) = a + b − c \\ Minus vor der Klammer: alle Vorzeichen drehen.  a − (b − c) = a − b + c \\ Zahl mal Klammer: jedes Glied malnehmen. \quad 3 · (x + 4) = 3x + 12}
